\documentclass[12pt]{article}
\usepackage{amsmath, amssymb}
\usepackage{geometry}
\usepackage{setspace}
\usepackage[utf8]{inputenc}
\usepackage{hyperref}
\usepackage{enumitem}
\usepackage{fancyhdr}
\usepackage{xcolor}

\geometry{margin=1in}
\setstretch{1.25}

\definecolor{citecolour}{RGB}{0, 0, 128}
\definecolor{forestgreen}{RGB}{34, 139, 34}
\hypersetup{
    colorlinks,
    linkcolor={citecolour},
    citecolor={citecolour},
    urlcolor={forestgreen}
}

\pagestyle{fancy}
\fancyhf{}
\lhead{Jacob Edenhofer}
\rhead{Research Design Proposals}
\cfoot{\thepage}
\setlength{\headheight}{15pt}

\title{Assignment -- Building Out Your Theory}
\author{Jacob Edenhofer\footnote{\href{mailto:jacob.edenhofer@nuffield.ox.ac.uk}{jacob.edenhofer@nuffield.ox.ac.uk}}}
\date{06/05/2026}

\begin{document}

\maketitle

\section*{Purpose}

The first assignment asked you to set up the theoretical argument at the centre of your research design proposal: the explanandum, contribution, type of theory, concepts, building blocks, central mechanism, closest rival, and scope. This assignment assumes that setup is already on the page. Do not simply repeat it.

The aim is not to design an empirical test. It is to begin formalising your informal setup. Formalising does not have to mean writing equations. It means turning prose into a stripped-down structure: who the actors are, what they choose, what they want, what they know, what is fixed, what is generated by the argument, and what changes when one part of the setup changes. Think of this as moving from a promising setup to the analytical core of the argument.

\section*{What to Submit}

Email a single document of approximately 1,000--1,500 words to me. You may use the section headings below, but feel free to merge sections if your argument reads more naturally that way. You do not need to use mathematics, though you may include notation, diagrams, a timeline, or a simple game tree if they clarify the argument.

\section*{Some Terms}

This assignment uses a few terms from model-building. Use them in a simple way.

\begin{description}[leftmargin=2.8cm, style=nextline]
\item[Target] The part of the real world you are trying to understand. For example: voters deciding whether to punish incumbents, parties screening candidates, firms lobbying regulators, or subnational governments bargaining with central governments.
\item[Model] A simplified representation of the target. A model can be formal, verbal, graphical, or computational. In this assignment, a model just means a disciplined simplified version of your argument.
\item[Formalise] To make the parts of an argument explicit enough that a reader can see the inputs, choices, constraints, assumptions, and outputs. Formalising may involve equations, but it can also involve a table, timeline, decision tree, or carefully specified prose.
\item[Primitive] An input to the argument: an actor, preference, belief, information structure, cost, benefit, constraint, resource, institution, timing assumption, or rule of the game. A primitive is something the theory takes as given while deriving its logic.
\item[Endogenous outcome] Something generated inside the argument: an action, choice, belief, strategy, equilibrium, allocation, or aggregate pattern. Endogenous outcomes are what the mechanism produces.
\item[Mechanism] The answer to a ``why'' question. It is the actor-level logic that connects primitives to endogenous outcomes.
\item[Comparative static] A claim about how an endogenous outcome changes when one primitive changes while the other primitives are held fixed.
\item[Nearby model] A modified version of your setup that changes one feature while leaving the rest as similar as possible. Nearby models help identify which feature is crucial for a result.
\item[Representational feature] A part of the model that is meant to correspond to something in the target.
\item[Simplifying assumption] An assumption that leaves out real-world detail so the main logic is easier to see. For example, treating two actors as if they each have one goal, even though real actors care about many things.
\item[Technical assumption] An assumption that mainly makes the argument easier to state or solve, rather than carrying the substantive mechanism. For example, using a one-dimensional policy space or a simple functional form.
\end{description}

\section{Start From Assignment 1}

Begin with a short bridge from your first assignment. In no more than one paragraph, remind the reader of your explanandum and central mechanism. Do not redefine all your concepts, restate the full rival explanation, or rewrite the whole setup.

Then identify one representational choice you made in Assignment 1:
\begin{enumerate}[label=(\alph*)]
\item one actor, situation, or mechanism in your setup that is meant to correspond to something in the real world;
\item one feature of the real world you deliberately left out;
\item why leaving it out helps isolate the mechanism for now.
\end{enumerate}

This section should be brief. Its purpose is to orient the reader, not to redo Assignment 1.

\section{From Prose to a Minimal Setup}

Reorganise your Week 1 building blocks into the parts of a minimal setup. The goal is to make your informal argument precise enough that someone else could see how it might become a simple model.

First, list the most important primitives. You do not need more than five. For each primitive, say what kind of object it is: actor, preference, belief, information, cost, benefit, constraint, institution, timing assumption, or rule.

Second, list the main endogenous outcomes. These are the things your mechanism explains or generates. They might include an actor's choice, a strategy, a belief, a pattern of behaviour, a distribution of outcomes, or an equilibrium prediction.

Then answer:
\begin{enumerate}[label=(\alph*)]
\item What is the actor's choice set: the options the actor can choose among?
\item What does the actor prefer, value, or try to avoid?
\item What does the actor know or believe when choosing?
\item Which primitive is most central to the result?
\item Which endogenous outcome is the main explanandum?
\item Are there intermediate outcomes that matter for the mechanism? For example, a belief may change first, then an action, then an aggregate outcome.
\end{enumerate}

\section{Crucial Feature and Nearby Model}

Rather than restating the whole mechanism, identify the crucial feature of the setup: the feature without which the result would disappear, weaken sharply, or change meaning.

Your answer should:
\begin{enumerate}[label=(\alph*)]
\item name the crucial feature;
\item explain where it appears in the mechanism from Assignment 1;
\item describe a nearby model that removes or changes only that feature;
\item say whether the original result survives, weakens, disappears, or reverses in the nearby model;
\item explain what this teaches the reader about why your mechanism works.
\end{enumerate}

For example, ask: if I removed this information asymmetry, changed this cost, eliminated this actor's constraint, or made this belief common knowledge, would the result still follow?

\section{Assumptions and What They Do}

Identify three assumptions your argument needs. At least one should be substantive rather than merely definitional.

For each assumption, answer:
\begin{enumerate}[label=(\alph*)]
\item what the assumption says;
\item whether it is representational, simplifying, or technical;
\item why the mechanism needs it;
\item what would change if the assumption were relaxed.
\end{enumerate}

An assumption is \emph{representational} if it is meant to correspond to something important in the target. It is \emph{simplifying} if it removes detail mainly to make the main logic easier to see. It is \emph{technical} if it helps you state the argument cleanly but is not itself the substantive mechanism.

\section{Implications and Comparative Statics}

Derive two or three implications from the mechanism. These should be theoretical implications, not an empirical strategy. At least one should be a comparative static.

For each implication, state:
\begin{enumerate}[label=(\alph*)]
\item the implication in plain language;
\item the primitive that changes, such as a cost, benefit, constraint, belief, signal, or institutional rule;
\item the endogenous outcome that changes;
\item the actor-level reason the implication follows;
\item whether the implication is common to your argument and the rival, or distinctive to your argument.
\end{enumerate}

The goal is to show what follows from the mechanism. ``More $X$ leads to more $Y$'' is a start, but the stronger version explains why the relevant actor changes behaviour when $X$ changes, holding the other primitives fixed.

\section{Rival Logic, Without Repeating the Rival}

Return to the closest rival from the first assignment, but do not restate it at length. The task here is to sharpen the analytical difference between your mechanism and the rival.

Your answer should include:
\begin{enumerate}[label=(\alph*)]
\item the rival mechanism in one sentence;
\item one implication from Assignment 1 that is common to your argument and the rival;
\item one implication that is more clearly distinctive after the work you have done in this assignment;
\item the assumption under which your mechanism is more plausible than the rival;
\item the assumption under which the rival might dominate your mechanism.
\end{enumerate}

This section is still theoretical. The point is not to say how you would test the rival. The point is to make clear why the two arguments are analytically distinct.

\section{Scope and Breakdown}

Do not provide a full scope statement again. Instead, identify the most important boundary condition revealed by the previous sections.

Your answer should include:
\begin{enumerate}[label=(\alph*)]
\item the representational feature that most limits where your argument travels;
\item the kind of setting where that feature is likely to be absent;
\item the assumption whose failure would most clearly break the mechanism;
\item the condition under which the result might reverse.
\end{enumerate}

\section*{A Detailed Worked Example: Delegation}

This example draws on a standard delegation model discussed in Bueno de Mesquita et al.'s \emph{Theory and Credibility}. The target is political delegation: why would a principal, such as a legislature, minister, or voter, allow an agent, such as a bureaucracy, agency, or representative, to choose policy on her behalf?

\paragraph{Target and representation.}
The real-world target is a setting in which one actor has authority but another actor has expertise. The model represents three features of that target. First, the principal cares about the final policy outcome. Second, the agent may know more about how policies translate into outcomes. Third, the principal and agent may not want exactly the same outcome. The model leaves out many real-world complications: multiple principals, legal procedures, lobbying, reputation, ideology across many dimensions, and repeated interaction. Those omissions are acceptable if the simplified setup isolates the trade-off between expertise and policy disagreement.

\paragraph{From prose to a minimal setup.}
The main primitives are:
\begin{enumerate}
\item A principal with ideal outcome $x_p$. This is the outcome the principal likes best.
\item An agent with ideal outcome $x_a$. This is the outcome the agent likes best.
\item A policy choice $\pi$.
\item An idiosyncratic shock $\varepsilon$, which captures uncertainty about how policy translates into outcomes.
\item A production rule $x=\pi+\varepsilon$, meaning that the final outcome depends on the chosen policy and the shock.
\item Information: the agent observes $\varepsilon$, while the principal knows only its distribution.
\end{enumerate}

These primitives are inputs. The model takes them as given in order to analyze the choice about delegation. Put in the language of formalisation: the principal's choice set is \{delegate, do not delegate\}; the principal wants an outcome close to $x_p$; the agent wants an outcome close to $x_a$; the agent knows $\varepsilon$; and the principal does not.

\paragraph{Endogenous outcomes.}
The main endogenous outcome is whether the principal delegates authority to the agent. Other endogenous outcomes include the policy chosen under delegation, the policy chosen without delegation, and the final outcome. If the principal delegates, the informed agent can choose $\pi=x_a-\varepsilon$, producing the agent's ideal outcome $x=x_a$. If the principal does not delegate, the principal chooses policy without observing $\varepsilon$, so she cannot perfectly control the final outcome.

\paragraph{Mechanism.}
The mechanism is an expertise-control trade-off:
\begin{enumerate}
\item The principal wants an outcome close to $x_p$ but is uncertain about how any policy $\pi$ will translate into the outcome $x$.
\item The agent observes the shock $\varepsilon$, so the agent can use policy to offset uncertainty.
\item Delegation therefore gives the principal informational benefits: the final outcome becomes more predictable.
\item But delegation also gives the agent control. If $x_a$ differs from $x_p$, the agent uses that control to implement the agent's preferred outcome.
\item The principal delegates when the informational benefit of using the agent's expertise is larger than the policy loss from letting an imperfect ally choose.
\end{enumerate}

\paragraph{Crucial feature and nearby model.}
The crucial feature is asymmetric information: the agent knows something about the policy environment that the principal does not know. To see this, consider a nearby model in which the principal also observes $\varepsilon$. Then the principal can choose $\pi=x_p-\varepsilon$ and obtain her own ideal outcome. Delegation no longer provides an informational advantage. The case for delegation collapses unless some other mechanism is added.

This nearby model clarifies why the delegation result is not simply about trust or hierarchy. The principal may dislike the agent's preferences and still delegate because the agent's information solves a problem the principal cannot solve alone.

\paragraph{Assumptions and what they do.}
One representational assumption is that the agent has relevant expertise. This corresponds to real settings in which bureaucrats, regulators, or local officials know more about implementation than the principal. If this assumption fails, delegation loses its central benefit.

A second representational assumption is that the principal and agent may have different ideal outcomes. This captures political conflict between the delegating actor and the actor receiving authority. If their preferences are identical, delegation becomes easier because the control cost disappears.

A simplifying assumption is that policy and outcomes are one-dimensional. Real policy often has many dimensions. The assumption helps isolate the expertise-control trade-off, but relaxing it could change the result if the agent is aligned with the principal on one dimension and misaligned on another.

A technical assumption is the additive production rule $x=\pi+\varepsilon$. This is not usually the substantive claim. It is a tractable way to represent uncertainty. If the mapping from policy to outcomes were more complex, the basic mechanism might survive, but the exact delegation condition could change.

\paragraph{Implications and comparative statics.}
The first implication is the ally principle: delegation becomes more attractive as the agent's ideal outcome moves closer to the principal's ideal outcome. The changing primitive is preference distance, $|x_a-x_p|$. The endogenous outcome is the principal's delegation decision. The actor-level logic is that an aligned agent imposes a smaller control cost.

The second implication is the uncertainty effect: delegation becomes more attractive as the principal faces more uncertainty about how policy translates into outcomes. The changing primitive is the variance or spread of $\varepsilon$. The endogenous outcome is again the delegation decision. The actor-level logic is that the agent's expertise becomes more valuable when uninformed policy choice is more likely to produce an outcome far from what the principal wants.

A third implication concerns breakdown. If the agent is very poorly aligned with the principal, delegation may fail even when the agent has expertise. The changing primitive is again preference distance, but now the implication is about the limit of the mechanism: expertise is not enough if the policy loss from agency control is too large.

\paragraph{Rival logic.}
A rival mechanism might explain delegation as capacity relief rather than expertise. On that account, the principal delegates because she lacks time, staff, or administrative capacity, not because the agent has superior information about the policy environment. Both mechanisms imply more delegation when tasks become more complex. They differ in why complexity matters. In the expertise mechanism, complexity matters because the agent can choose better policy conditional on information. In the capacity-relief mechanism, complexity matters because the principal cannot process or implement all tasks herself.

The expertise mechanism is more plausible when the agent has information that changes which policy should be chosen. The capacity-relief mechanism is more plausible when tasks are numerous or costly but do not require much specialized information.

\paragraph{Scope and breakdown.}
The delegation model should travel to settings where principals value outcomes, agents have relevant expertise, and policy authority can be transferred. It should not apply well where the principal can cheaply acquire the same information, where agents cannot act on their information, or where delegation is legally impossible. The result might reverse if delegation creates severe agency losses, for example when the agent's preferences are so far from the principal's that expertise no longer compensates for loss of control.

\section*{Checklist}

Before submitting, ask whether your memo lets a reader answer the following questions:
\begin{itemize}
\item Have I avoided simply rewriting Assignment 1?
\item What representational choice from Assignment 1 matters most?
\item Have I translated my prose into a minimal setup with actors, choices, preferences, information, and constraints?
\item Which parts of the setup are primitives?
\item Which outcomes are generated inside the argument?
\item What is the crucial feature without which the result would change?
\item Which assumptions are doing the most analytical work?
\item What changes when a key primitive increases or decreases?
\item What is sharper now about the difference between my mechanism and the rival?
\item Which boundary condition follows from the theory-building work above?
\end{itemize}

\end{document}
